%%
%% frontmatter/abstract_zh.tex -- Chinese (zh-TW) abstract + keywords
%% Exact text locked in the quick-260713-mkh CONTEXT; numbers byte-consistent
%% with the English abstract. bkai (標楷體-like) via CJKutf8.
%%
\begin{CJK}{UTF8}{bkai}
\chapter*{摘要}
\addcontentsline{toc}{chapter}{\texorpdfstring{\protect\begin{CJK}{UTF8}{bkai}摘要\protect\end{CJK}}{Chinese Abstract}}

弱監督影片異常偵測僅以影片層級標籤定位異常事件，然而現有方法多仰賴單一視覺模態，且假設部署環境固定不變。本論文提出一套雙模態融合框架，結合凍結 CTR-GCN 編碼器之骨架動態特徵與凍結視覺語言骨幹網路之視覺特徵，以多實例學習排序損失訓練，僅學習輕量之融合頭。研究涵蓋四種視覺骨幹網路（CLIP ViT-B/16、SigLIP2 ViT-B/16、SigLIP2 SO400M、SigLIP2 Giant），於 UCF-Crime 與 XD-Violence 基準資料集上進行評估。

閘控融合以學習式模態加權於 UCF-Crime 達到 82.5\% 幀級 AUC（SigLIP2 Giant），於 XD-Violence 達到 78.7\% AP（SigLIP2 SO400M），皆為三種子平均。骨架特徵帶來小而一致的互補增益：閘控融合於全部八組骨幹—資料集設定中皆不低於純視覺基線，其中六組嚴格優於基線。上述結果落後於領域中經微調與文字對齊之領先方法（約 88--91\% UCF-Crime AUC），此差距源於本論文刻意採用之凍結、無文字分支設計。本論文之貢獻為可重現之單 GPU 雙模態流程與面對影像劣化之免標籤強健性研究，而非新的最高分數。

本論文另探討測試時期自適應（TTA）：於 20 種合成劣化條件、三種子下，以熵最小化為基礎之兩種代表性 TTA 方法（TENT 與 SAR）受架構限制僅能調整融合頭之 LayerNorm 參數，AUC 變化均小於 0.1 個百分點，為論文中深入分析之結構性零結果。所提出之免標籤、免調參的判別可靠度重加權方法，於視覺語言串流之判別訊號崩潰時，將閘控融合導向對劣化更穩健之骨架串流：於全部四種骨幹之劣化輸入 AUC 均有提升（平均 +1.2 點，於單一視覺串流崩潰最嚴重之條件最高 +13.2 點），採轉導式協定且不宣稱於乾淨資料上有增益。

\vspace{1.5em}
\noindent\textbf{關鍵字：}影片異常偵測、弱監督學習、多模態融合、骨架動作表徵、視覺語言模型、測試時期自適應
%% zh-TW abstract wording approved by the author 2026-07-14.
%% Post-defense amendment 2026-07-29 (author-approved in session): 「免標籤之強健性研究」
%% → 「面對影像劣化之免標籤強健性研究」, resolving committee question on 強健性 scope;
%% EN counterpart "corruption-robustness" applied the same day.
\end{CJK}
